\documentclass[11pt,a4paper]{article}
\usepackage[utf8]{inputenc}
\usepackage[top=0.5in, bottom=1in, left=1in, right=1in]{geometry}
\usepackage{url}
\usepackage{graphicx}
\usepackage{parskip}
\usepackage[usenames,dvipsnames]{xcolor}
\usepackage{titlesec}
\usepackage{ebgaramond}
\usepackage{tabularx}
\usepackage{enumitem}
\usepackage{multirow}
\usepackage{fix-cm}
\usepackage{array}
\usepackage{tikz}

%\usepackage{charter}
%\usepackage{times}

\makeatletter
\newcommand\CVsize{\@setfontsize\Huge{60}{60}}
\makeatother  

% Link formatting
\usepackage{hyperref}
\definecolor{linkcolour}{rgb}{0,0.2,0.6}
\hypersetup{colorlinks,breaklinks,urlcolor=linkcolour, linkcolor=linkcolour}

%Section formatting
\titleformat{\section}{\Large\scshape\raggedright}{}{0em}{}[\titlerule]
\titlespacing{\section}{0pt}{3pt}{3pt}

%\newcommand{\sbt}{\,\begin{picture}(-1,1)(-1,-3)\circle*{3}\end{picture}\ }

\begin{document}
\pagestyle{empty}

%--------------------TITLE-------------

\hspace{-7pt}\begin{tikzpicture}
\node[inner sep=0pt] (l) at (0,0) {};
\node[inner sep=0pt] (m) at (0.5\linewidth,0) {};
\node[inner sep=0pt] (r) at (1.008\linewidth,0) {};

\node[inner sep=0pt, anchor=west] at (l) {\CVsize CV};

\node at (m) {\parbox{\linewidth}{\centering \Huge Peter Kvillegård \\ \LARGE Fysioterapeut\\[5pt] \normalsize Tel: 0760330019 \hspace{0.3em} \textbar \hspace{0.3em} E-post: peterkvillegard@gmail.com}};

\node[inner sep=0pt, anchor=east] (pic) at (r)
{\includegraphics[width=.16\linewidth]{portrait}};
\end{tikzpicture}

%--------------------SECTIONS-----------------------------------
%Section: Utbildning
\vspace{-5pt}
\section{Utbildning}
\hspace{-4pt}
\begin{tabularx}{\textwidth}{r|X}	
	\multirow{2}{*}{
		\begin{tabular}{@{}r@{}}
			Okt 2017\\
			-- Juni 2020\\
			\footnotesize{(Planerad}\\
			\footnotesize{examen)}
		\end{tabular}
	}
	& Kandidatexamen i Fysioterapi, Lunds universitet, Lund\\
	&\footnotesize{\vspace{-5pt}
		 \begin{itemize}[leftmargin=10pt, topsep=-12.5pt]

		\item Praktikperiod 1: Geriatrikavdelning 2, Malmö. Här lärde jag mig att bemöta och hjälpa patienter med olika behov, och att dokumentera i Melior.
		\item Praktikperiod 2: Barn- och ungdomshabiliteringen, Lund. Här lärde jag mig mer om patienter med olika neuro\-logiska tillstånd, samt undersökning och rehabiliteringsprinciper.
		\item Kandidatuppsats: ``EMG-analys av m. deltoideus och m. infraspinatus vid utåt\-rotations\-övningar med och utan adduktionskraft.''
		\end{itemize}\vspace{-30pt} 
	}\\
	\multicolumn{2}{c}{} \\
\end{tabularx}

%Section: Övriga utbildningar
\section{Övriga utbildningar}
\begin{tabularx}{\textwidth}{r|X}
	
	Okt. 2016& Certifierad massageterapeut, Axelsons gymnastiska institut, Göteborg\\
	-- Mars 2017& \footnotesize{Utbildningen innefattade bl.a. klassisk svensk massage, taktil massage, och barnmassage.} \\
	\multicolumn{2}{c}{} \\
	
	Aug. 2005& Kurser på Ölands folkhögskola, Skogsby \\
	--Maj 2006&\footnotesize{\vspace{-5pt}
		\begin{itemize}[leftmargin=10pt, topsep=-12.5pt]
			\item Hälsa och livskvalitet
			\item Skrivarskolan
		\end{itemize}\vspace{-30pt} 
	}\\
	\multicolumn{2}{c}{} \\
	
	Sep. 2004& Kurser på Malmö högskola, Malmö\\
	-- Dec. 2004&\footnotesize{\vspace{-5pt}
		\begin{itemize}[leftmargin=10pt, topsep=-12.5pt]
			\item Hälsa och samhälle, 7.5 hp
			\item Medicinsk vetenskap, 7.5 hp
		\end{itemize}\vspace{-30pt} 
	}\\
	\multicolumn{2}{c}{} \\	
		
	Jan. 2003& Kurs på Linnéuniversitetet, Kalmar\\
	-- Mars 2003&\footnotesize{\vspace{-5pt}
		\begin{itemize}[leftmargin=10pt, topsep=-12.5pt]
			\item English for science and technology, 7.5 hp
		\end{itemize}\vspace{-30pt} 
	}\\
	
	\multicolumn{2}{c}{} \\
	
	Aug. 1998& Gymnasium, Naturvetenskapsprogrammet teknisk gren, Mönsterås gymnasium,\\
	-- Juni 2001&Mönsterås\\
	\multicolumn{2}{c}{} \\
\end{tabularx}

%Section: Yrkeslivserfarenhet
\section{Yrkeslivserfarenhet}
\begin{tabularx}{\textwidth}{r|X}	
	Mars 2017 & Massageterapeut, Högsby \\
	-- Sept 2017 &\footnotesize{Över sommaren masserade jag kunder vilket utvecklade min palpationsfärdighet, och jag träffade många olika människor med olika besvär och behov av hjälp.}\\
	\multicolumn{2}{c}{} \\

	Feb 2007 & Maskinoperatör/programmerare på Stranda AB, Södra Bäckebo \\
	-- Sept 2016&\footnotesize{På denna arbetsplats fick jag många års erfarenhet av vad det innebär att jobba i ett team, att kunna kommunicera och ha ansvar. Att genom samarbete lösa problem.}\\
	\multicolumn{2}{c}{} \\
	
	Okt 2001& Massaläggare på Scania AB, Oskarshamn \\
	-- Mars 2004&\footnotesize{Mitt första jobb lärde mig god arbetsmoral och vikten av en organiserad arbetsplats.}\\
	\multicolumn{2}{c}{} \\

\end{tabularx}

%Section: Övriga meriter
\section{Övriga meriter}
\begin{tabularx}{\textwidth}{X}
	\vspace{-7pt}
\begin{itemize}[leftmargin=0.8em]
%	\item ?
%	\item ?
	\item Erfarenhet av olika typer av yoga samt aerobisk och muskelstärkande träning.
\end{itemize}
\end{tabularx}

%Section: Styrkor
\section{Styrkor}
\begin{tabularx}{\textwidth}{X}
	\vspace{-7pt}
	\begin{itemize}[leftmargin=0.8em]
		\item Bra på att bemöta människor
		\item Noggrann och tålmodig
		\item Stort intresse för att ta till mig ny kunskap
	\end{itemize}
\end{tabularx}

\section{Intressen och aktiviteter}
\begin{tabularx}{\textwidth}{X}
Gå omkring i skogen, läsa alla typer av böcker, programmering, besöka British Columbia med min flickvän, vegetarisk matlagning, spela gitarr, och ta hand om djur.\\
\\
\end{tabularx}


\section{Referenser}
\begin{tabularx}{\textwidth}{X}
Jeannette Unge, handledare för kandidatuppsats,\\ Tel: +46 46 222 89 58, E-post: jeannette.unge@med.lu.se
\\
\end{tabularx}

\end{document}